\subsection{Suposiciones}

Para el diseño de la base de datos se han considerado ciertas suposiciones que permiten definir su estructura y funcionamiento. Se asume que el sistema será implementado sobre un  gestor de bases de datos relacional (SQLite), con soporte para claves primarias, foráneas y restricciones de integridad.

Asimismo, se supone que el sistema será utilizado a través de una aplicación externa (de escritorio), encargada de la carga, consulta y visualización de los datos. En este sentido, se espera que los usuarios finales (entrenadores, organizadores o administradores) no interactúen directamente con la base de datos, sino mediante interfaces diseñadas para tal fin.

Se considera además que el volumen de datos será moderado, correspondiente a competencias, clubes y jugadores en el ámbito formativo, por lo que no se requieren inicialmente mecanismos avanzados de distribución o procesamiento masivo de datos. Sin embargo, el diseño contempla la posibilidad de crecimiento futuro.

\subsection{Restricciones}

El diseño de la base de datos se encuentra condicionado por diversas restricciones que impactan en su estructura. En primer lugar, se ha optado por un modelo relacional, lo cual implica la necesidad de respetar principios de normalización para evitar redundancias y anomalías en los datos.

Otra restricción relevante es la necesidad de mantener integridad entre entidades, lo que implica la correcta definición de claves primarias y foráneas, especialmente en relaciones como inscripciones o asociaciones entre clubes y competencias.

También se consideran limitaciones propias del entorno de implementación, como la capacidad de almacenamiento disponible, el rendimiento del sistema y la simplicidad requerida para su mantenimiento. Estas restricciones llevan a evitar diseños excesivamente complejos o altamente desnormalizados.

Por último, pueden existir restricciones funcionales derivadas de los requerimientos del sistema, como la unicidad de ciertos datos (por ejemplo, identificadores de jugadores o competencias) y la obligatoriedad de ciertos atributos para garantizar la consistencia de la información.

\subsection{Riesgos}

El diseño de la base de datos presenta ciertos riesgos que deben ser considerados. Uno de los principales es la posible inconsistencia de los datos debido a errores en la carga de información, especialmente si no se aplican validaciones adecuadas desde la aplicación que interactúa con la base de datos.

Otro riesgo es la evolución de los requerimientos del sistema, lo cual podría implicar modificaciones en el modelo de datos. Si el diseño no es lo suficientemente flexible, estos cambios podrían generar costos adicionales de mantenimiento o migración.

Asimismo, existe el riesgo de bajo rendimiento en consultas a medida que crece el volumen de datos, particularmente si no se definen adecuadamente índices o estructuras de optimización.

Para mitigar estos riesgos, se proponen diversas estrategias: implementar restricciones de integridad a nivel de base de datos, validar los datos desde la aplicación, diseñar el modelo con criterios de escalabilidad y mantener una documentación actualizada que facilite futuras modificaciones. Además, se recomienda realizar pruebas periódicas de rendimiento y consistencia de los datos.
